\section{プラズマ状態、電子衝突および準位収支}
\subsection{領域ごとの入力状態と計算対象}
領域$z$の状態を、電子密度$\nelec(z)$、エネルギー空間の電子確率密度$\fe(E;z)$、外部供給密度$N_s(z)$、解く励起準位密度$n_i^*(z)$で表す。ガス圧力$p$とガス温度$T_g$から親ガスの密度を
\begin{equation}
N_{\mathrm{gas}}=\frac{p}{k_{\mathrm B}T_g},\qquad N_s=x_sN_{\mathrm{gas}}
\label{eq:gas}
\end{equation}
とする。$x_s$は入力のモル分率である。親ガスは圧力・温度・分率から、ラジカル、準安定状態、残留ガスは領域別の密度指定から構成する。例えばCF$_4$/O$_2$/Arでは、親ガス組成に加えてF、CF、CF$_2$、Oの密度を与える。この入力密度を測定値または外部の化学モデルから定めることで、対象とする励起動力学に計算を集中できる。

各密度は一つの入力欄が所有する。電子、親ガス、ラジカル、準安定状態と、CRで解く準位の対応を検査し、同一種の二重指定は入力エラーとする。\code{solve: false}の準位は既知の密度として扱い、その準位からの発光と、解く下準位への放射供給を計算できる。これにより、独立に測定した準安定・励起準位密度を発光計算へ組み込める。

\subsection{電子エネルギー分布関数と温度の意味}
本コードは、電子エネルギー$E$を変数とする確率密度$\fe(E)$を用いる。$\fe(E)\,dE$は、電子が$E$から$E+dE$の区間に存在する割合である。速度空間分布や電子エネルギー確率関数（EEPF）の資料を使う場合は、この定義へ換算した分布を入力する。
\begin{equation}
\int_0^\infty \fe(E)\,dE=1,\qquad [\fe]=\mathrm{eV}^{-1},\qquad
\langle E\rangle=\int_0^\infty E\fe(E)\,dE .
\label{eq:pdf}
\end{equation}
Maxwell分布の入力$\Te$は、$k_{\mathrm B}T$をeVで表した量であり、十分な積分範囲では
\begin{equation}
f_{\mathrm M}(E;\Te)=\frac{2}{\sqrt\pi}\frac{\sqrt E}{\Te^{3/2}}
\exp\!\left(-\frac{E}{\Te}\right),\qquad \langle E\rangle=\frac32\Te .
\label{eq:maxwell}
\end{equation}
励起閾値より上の分布が速度係数を左右するため、平均エネルギーが同じでも、高エネルギー成分が異なれば発光は異なり得る\cite{hagelaar}。

Druyvesteyn型は、実装の\code{te\_eV}を尺度$E_D$と読み替えて
\begin{equation}
f_{\mathrm D}(E;E_D)=\frac{2\sqrt E}{E_D^{3/2}\Gamma(3/4)}
\exp\!\left[-\left(\frac E{E_D}\right)^2\right],\qquad
\langle E\rangle=E_D\frac{\Gamma(5/4)}{\Gamma(3/4)}
\label{eq:druy}
\end{equation}
と書ける。式\eqref{eq:druy}の係数と平均値は無限区間の解析値であり、コードは有限格子で数値的に正規化する。$E_D$は分布の尺度である。Maxwell分布との比較には、両分布の平均エネルギーを式\eqref{eq:maxwell}、\eqref{eq:druy}から求める。

二温度Maxwell分布は
\begin{equation}
f_{\mathrm{bi}}(E)=(1-\alpha)f_{\mathrm M}(E;T_c)+\alpha f_{\mathrm M}(E;T_h),
\qquad 0\leq\alpha\leq1,\quad T_h\geq T_c>0
\label{eq:bimax}
\end{equation}
である。$\alpha$は高温成分に属する電子の粒子割合で、$T_c$、$T_h$とともに指定する。二温度表現により、少数の高エネルギー電子が励起率へ及ぼす影響を調べられる。外部で得た分布を直接用いる場合は表形式EEDFを選び、内部格子への線形補間、表の範囲外のゼロ補完、再正規化を経て、速度係数と発光を順解析する。

数値計算では、実際に使う格子$[E_{\min},E_{\max}]$上で各分布を再正規化する。積分範囲は励起に寄与する高エネルギー側まで含める。上端付近の確率質量と分布振幅を診断し、必要に応じて$E_{\max}$を拡大して速度係数の収束を確認することで、裾の切断が発光へ及ぼす影響を評価する。

\subsection{電子衝突速度係数と断面積}
反応$r$の断面積$\sigma_r(E)$とEEDFから、速度係数を
\begin{equation}
k_r[\fe]=\int_0^\infty \sigma_r(E)\sqrt{\frac{2eE}{m_{\mathrm e}}}\,\fe(E)\,dE,
\qquad [k_r]=\mathrm{m^3\,s^{-1}}
\label{eq:rate}
\end{equation}
で得る。$e$はeVをJへ換算する素電荷、$m_{\mathrm e}$は電子質量である。エネルギーの単位換算は速度$\sqrt{2eE/m_{\mathrm e}}$に含まれ、$\fe(E)\,dE$は無次元の確率を与える。断面積を共有エネルギー格子へ線形補間し、台形則で積分する。一定の速度係数や、閾値を持つ解析的な代用断面積も選べるため、解析解との比較や機構の感度調査を簡潔に構成できる。実ガスの定量計算では、対象過程に対応する断面積と出典を選定する。

断面積CSVは\code{energy\_eV,sigma\_m2}の二列とし、エネルギーの単調増加、有限性、断面積の非負性を検査する。積分処理では表の外をゼロで補完し、その範囲のEEDF確率を未被覆量として記録する。利用する電子分布の裾に対応したエネルギー範囲の表を用意することが、率計算の入力条件となる。物理閾値から始まる表では、先頭行の断面積をゼロとして\code{zero\_below\_threshold}を宣言すると、閾値未満を既知のゼロ領域として被覆に含める。

断面積の出典、過程名、エネルギー範囲、SHA-256を記録するため、LXCat等から抽出した入力を追跡できる\cite{lxcat}。被覆率は電子確率分布と表のエネルギー範囲の対応を示す。速度係数の精度を調べるときは、被覆率に加え、積分核$\sigma v\fe$、励起閾値、および積分範囲の拡大に対する率の変化を確認する。

\subsection{定常線形CR方程式の組立て}
解く準位$i$への外部供給を$S_i$、追跡する準位$j$から$i$への一次移行頻度を$\nu_{j\to i}$、追跡外への損失を$\nu_{i\to\mathrm{out}}$とする。
\begin{equation}
\frac{dn_i^*}{dt}=S_i+\sum_{j\ne i}\nu_{j\to i}n_j^*
-\left(\sum_{j\ne i}\nu_{i\to j}+\nu_{i\to\mathrm{out}}\right)n_i^* .
\label{eq:balance}
\end{equation}
電子分布と衝突相手密度を固定すれば、すべての$\nu$は既知のs$^{-1}$である。定常化により
\begin{equation}
M\boldsymbol n^*=\boldsymbol b,\qquad
M_{ii}=\sum_{j\ne i}\nu_{i\to j}+\nu_{i\to\mathrm{out}},\qquad
M_{ij}=-\nu_{j\to i}\ (i\ne j),\qquad b_i=S_i
\label{eq:matrix}
\end{equation}
を解く。外部密度$N_s$からの供給は$S_i\mathrel{+}=\nu_{s\to i}N_s$として右辺へ入る。未知準位を出発点とする過程は、対角の損失と、追跡する到達準位への非対角供給を同時に加える。各過程の行列寄与と生成・損失収支を同じ記録から構築するため、特定の項が結果を支配する理由を追跡できる。

\begin{table}[htbp]
\centering\small
\begin{tabularx}{\linewidth}{p{31mm}p{38mm}p{24mm}Y}
\toprule
過程 & 有効移行頻度$\nu$ & 係数の単位 & 代表的用途 \\
\midrule
電子衝突 & $\nelec k_e[\fe]$ & m$^3$ s$^{-1}$ & 基底・準安定からの励起、宣言した準位移行 \\
一次過程 & $k_1$ & s$^{-1}$ & 非放射損失、指定した一次移行 \\
二体過程 & $k_2N_M$ & m$^3$ s$^{-1}$ & 中性粒子による消光 \\
三体過程 & $k_3N_MN_N$ & m$^6$ s$^{-1}$ & 外部衝突相手を伴う有効移行・損失 \\
放射過程 & $\Aeff_{u\ell}$ & s$^{-1}$ & 放射損失、追跡下準位へのcascade \\
壁損失 & $\nu_{w,i}$ & s$^{-1}$ & 指定した準位の有効境界損失 \\
\bottomrule
\end{tabularx}
\caption{実装済みの移行・損失項。二体・三体過程では、衝突相手の外部密度を既知係数として線形性を保つ。未知準位を衝突相手に指定する入力は反応式の検査でエラーとする。}
\label{tab:terms}
\end{table}

\subsection{二準位の例：式とコードの対応}
外部基底粒子から準位1、2へ$S_1,S_2$が供給され、準位2から1へ放射cascade $A_{21}$があるとする。$D_i$を各準位の全損失頻度とすると
\begin{equation}
\begin{pmatrix}D_1&-A_{21}\\0&D_2\end{pmatrix}
\begin{pmatrix}n_1^*\\n_2^*\end{pmatrix}
=\begin{pmatrix}S_1\\S_2\end{pmatrix},\qquad
n_2^*=\frac{S_2}{D_2},\quad
n_1^*=\frac{S_1+A_{21}S_2/D_2}{D_1}.
\label{eq:twolevel}
\end{equation}
$D_2$には$A_{21}$自体も含む。消光を加えれば該当する$D_i$が増え、cascadeの除外は上準位の損失と下準位の供給の双方を変える。単一準位の場合は
\begin{equation}
n_u^*=\frac{\nelec N_s k_{s\to u}}
{\sum_\ell\Aeff_{u\ell}+\sum_M k_{q,uM}N_M+\nu_{w,u}+\nu_{\mathrm{other},u}}
\label{eq:single}
\end{equation}
となり、密度と温度以外の入力が発光へどう入るかを明瞭に読める。

この準定常解は、励起系の緩和が入力状態の変化に十分速く追従する条件で用いる。緩和時間は単独準位で概ね$D_u^{-1}$、結合系で行列の緩和固有モードを目安に評価する。利用者はこれを電子状態や供給密度の変動時間と比較し、解析する時間区間を選ぶ。時間平均スペクトルへ適用する場合は、発光の入力依存性に由来する平均化の影響も評価する。時間発展と輸送の連成は第\ref{sec:outlook}節の拡張課題で扱う。

\subsection{放射捕獲、分岐、壁損失}
放射捕獲は、指定した光学的厚さ$\tau$に対する簡略slab脱出係数で表す。
\begin{equation}
\beta(\tau)=\frac{1-\exp(-\tau)}{\tau},\qquad \beta(0)=1,\qquad
\Aeff_{u\ell}=\beta_{u\ell}A_{u\ell}.
\label{eq:escape}
\end{equation}
この実装は、指定した光学的厚さ$\tau$または脱出係数$0\leq\beta\leq1$から、有効放射率$\Aeff$を求めるslab近似である。同じ$\Aeff$を準位収支と発光の双方へ用いることで、捕獲による寿命の変化と外へ逃げる光子の生成を一貫して扱う。数値的には$\tau\leq10^{-12}$で1とし、slabの$\beta$を$[10^{-6},1]$に制限する。例えば放射損失だけを持つ単一経路では、$\beta$の減少が占有率の増加を伴い、定常状態の全脱出光子数が保たれる極限が得られる。光学的厚さの値は、対象線の吸収と観測条件に基づいて与える。

原子データは波長と放射率の集合として管理し、各ケースの衝突過程と組み合わせる。例えばArの2p1は667.7、750.4 nm、2p6は763.5、800.6、922.4 nmへの分岐を持つ。観測範囲外も含む分岐集合を$D_u$へ入れることで、上準位寿命と観測線の光子収率を求める。NISTの波長と$A$値を固定した種別packは、この集合を複数のケースで再利用するために有用である\cite{nist}。

壁損失は
\begin{equation}
\nu_{w,i}=\gamma_i f_{\mathrm{flux}}\frac{\bar v_i}{L},\qquad
\bar v_i=\sqrt{\frac{8k_{\mathrm B}T_g}{\pi m_i}},\qquad
L=\frac VA=\frac{RH}{2(H+R)}
\label{eq:wall}
\end{equation}
とする。$R,H$は円柱容器の半径・高さで、$L$の直接指定も可能である。$\gamma_i$は状態別または種別の損失確率、$m_i$は質量である。この項は、指定状態に対する境界損失を一次の消滅頻度として準位収支へ加える。\code{thermal\_flux\_factor}の既定値は互換性のため1であり、等方的な平均速さと通常の熱流束$N\bar v/4$を使う定義では$1/4$を明示する。非ゼロの壁損失を用いる状態には質量を指定し、$\gamma_i$と$L$を対象の壁面条件に合わせて定める。

\subsection{準位解と格子収束の確認}
準位系には密行列の直接解法を用いる。各未知準位を放射・衝突・壁損失で釣り合わせ、一意な定常解を持つ反応集合を与える。特異系は通常の予測時にもエラーとし、利用者が損失経路や固定密度を点検できるようにする。分子帯だけを計算するケースでは、空のCR系を使用できる。監査では行列のrank・条件数、相対残差$\|M\boldsymbol n^*-\boldsymbol b\|_2/\|\boldsymbol b\|_2$、および負の準位密度量を記録する。実装は負値を大きさによらずゼロへ置き換えるため、置換前の解と残差を用いて計算の受理を判断する。

格子収束の確認では、エネルギー格子と波長格子を個別に細分化し、同じ検出画素上で予測の相対L2差を比較する。この手順により、速度係数積分と線形状・画素積分に必要な分解能を決められる。反応集合と断面積の妥当性は、生成・損失収支の比較および第\ref{sec:benchmarks}--\ref{sec:outlook}節の検証手順で評価する。
